\section{The crux in special linear settings}\label{sec:crux}

The crux asks for a deficit below rate~$1$, and its undirected core asks whether storage codes on locally bipartite side
information of bounded degree can beat rate~$\frac12$~\cite[Question~6.1]{FourthPaper2026}. For all codes this is open beyond
the classes of the companion paper, \emph{Nested Elimination and Half-Rate Bounds for Storage Codes}~\cite{SeventhPaper2026},
which proves the half bound for every storage code on triangle-free graphs of maximum degree three, among other classes, with
the guessing number equal to the fractional matching number there. This section proves two linear cases in which a natural
source of high-rate codes fails: codes lifted along towers of covers in characteristic~$2$ (\cref{sec:crux-towers}), and scalar
codes at degree four whose exceptional vertices are few or far apart (\cref{sec:crux-deg4}). By \cref{sec:linear}, linear codes
are the relevant ones for linear machines.

\subsection{Lifted codes along towers of covers}\label{sec:crux-towers}

Let $H$ be a finite connected graph with vertex set~$V$, $\card V=N$, and let $M\in\F_2^{V\times V}$ have unit diagonal and
off-diagonal support in the edges of~$H$; the entries $M_{uv}$ and $M_{vu}$ may differ. Its kernel is a linear storage code
on~$H$~\cite[Lemma~6.14]{FourthPaper2026}. Let $\mathcal T$ be the universal cover of~$H$, a tree, and $\pi=\pi_1(H)$, a free
group acting freely on~$\mathcal T$. A \emph{finite cover} is $\mathcal T/\Lambda$ for a subgroup $\Lambda\le\pi$ of finite
index; more generally, a finite $\pi$-set~$X$ defines a cover with vertex set $X\times V$. Choosing a spanning tree and a base
vertex turns~$M$ into a matrix $\tilde M\in\mathrm{Mat}_V(\F_2[\pi])$, whose entry $(u,v)$ is the sum of $M_{uv}g_e$ over the
edges~$e$ from $u$ to~$v$, where $g_e\in\pi$ is the closed walk along the tree to~$u$, across~$e$ and back along the tree, and
$\tilde M_{vv}=1$. The \emph{lift} $M_X$ is the image of $\tilde M$ under the permutation representation of~$\pi$ on
$\F_2[X]$, applied entrywise: a matrix with unit diagonal and support in the edges of the cover, whose kernel is a linear
storage code there, of rate $\dim\ker M_X/(N\card X)$. The \emph{inner rank} $\irk\tilde M$ is the least~$\rho$ such that
$\tilde M=PQ$ with $P\in\mathrm{Mat}_{N\times\rho}(\F_2[\pi])$ and $Q\in\mathrm{Mat}_{\rho\times N}(\F_2[\pi])$. The group algebra
of a free group is a free ideal ring, and its inner rank equals the rank over its universal field of fractions, Cohn's free
field~\cite{Cohn2006}. The \emph{universal rate} is
\[
  r_U(M)=\bigl(N-\irk\tilde M\bigr)/N .
\]
Let $\chi\colon\pi\to\Z/2$ be the parity of the length of closed walks, $\pi^+=\ker\chi$, $\hat\pi$ the pro-$2$ completion of
$\pi$ and $\hat\pi^+$ the closure of~$\pi^+$ in it. A cover $\mathcal T/\Lambda$ is bipartite if and only if $\Lambda\le\pi^+$.

\begin{theorem}[The equivariant crux in characteristic two]\label{thm:eqcrux}
\partformal{TwistPairing.lean}{lift2, twistPairing, finrank_ker_lift2, excess_lift2, rate_lift2_le}{parts (a), (b), (d),
(e) and~(f), and in~(c) the passage from the tower to its levels: the Lean theorems are the step from one level to the next,
for a double cover given by its twisted pairs.}
Let $H$ be non-bipartite.
\begin{enumerate}[label=(\alph*)]
\item Every finite cover~$X$ satisfies $\operatorname{rate}(M_X)\ge r_U(M)$.
\item $r_U(M)\le\frac12$.
\item Along every tower $\pi=\Lambda_0\ge\Lambda_1\ge\cdots$ in which each $\Lambda_{n+1}$ has index~$2$ in~$\Lambda_n$, the rate
is non-increasing. More precisely, with $N_n=2^nN$ and $M_n$ the lift at level~$n$, the excess $e_n=\dim\ker M_n-N_n/2$
satisfies $e_{n+1}=2e_n-r_n$, where $r_n\ge0$ is the rank of the \emph{twist pairing}
$\ker M_n^\T\times\ker M_n\to\F_2$, $(y,x)\mapsto\sum(M_{uv}y_ux_v+M_{vu}y_vx_u)$, the sum over the edges $uv$ of level~$n$
that the double cover twists.
\item Let $C=\bigcap_n\overline{\Lambda_n}\le\hat\pi$. Some level of the tower is bipartite if and only if $C\le\hat\pi^+$, and
from that level on the rate is at most~$\frac12$.
\item Let $\iota\colon\pi\to\mathbf F$ be an injective homomorphism into a free-by-$\Z_2$ pro-$2$ group~$\mathbf F$, where
$\Z_2$ denotes the $2$-adic integers (for instance a finitely generated free pro-$2$ group), and let
$\mathbf F=\mathbf N_1\ge\mathbf N_2\ge\cdots$ be open normal subgroups with trivial intersection. Then
$\operatorname{rate}(M_{\iota^{-1}(\mathbf N_j)})\to r_U(M)$.
\item For the Clebsch graph and $M=I+A$, $r_U(M)=0$.
\end{enumerate}
\end{theorem}

\begin{outline}[Outline of the formal proof of the step in~(c)]
In Lean the double cover of level~$n$ has vertex set $V_n\times\{0,1\}$ and is given by its twisted pairs; the lift
(\lean{lift2}) connects equal sheets by the untwisted entries of~$M_n$ and opposite sheets by the twisted ones, whose matrix
is $M_n^{\mathrm{tw}}$ (\lean{twistPart}). The field is any field of characteristic~$2$. In the coordinates $a=x_0+x_1$,
$b=x_1$ (\lean{sheetEquiv}) the lift becomes $(a,b)\mapsto(M_na,M_nb+M_n^{\mathrm{tw}}a)$ (\lean{blockMap_sheetEquiv}).
(Equivalently, in $\F_2[\Z/2]\cong\F_2[t]/(t^2)$ with $g=1+t$, the lift is $\varphi=M_n+tM_n^{\mathrm{tw}}$ in the basis
$1,t$; the Lean proof works with the block map.) Its kernel consists of the
pairs $(a,b)$ with $M_na=0$ and $M_nb=M_n^{\mathrm{tw}}a$. Since $\operatorname{im}M_n=(\ker M_n^\T)^\perp$
(\lean{mem_range_iff_dotProduct_ker_transpose}), the admissible~$a$ form the kernel of the twist pairing
$x\mapsto(y\mapsto y^\T M_n^{\mathrm{tw}}x)$ on $\ker M_n$ (\lean{twistPairing}), of dimension $\dim\ker M_n-r_n$, and each
admits a coset of $\ker M_n$ of choices of~$b$ (\lean{finrank_ker_blockMap}). So $\dim\ker M_{n+1}=2\dim\ker M_n-r_n$
(\lean{finrank_ker_lift2}), which is the recursion $e_{n+1}=2e_n-r_n$ (\lean{excess_lift2}), and the rate does not increase
(\lean{rate_lift2_le}). In Lean the twisted pairs are an arbitrary relation; for a double cover they are the two
orientations of the twisted edges, and then $y^\T M_n^{\mathrm{tw}}x$ is the sum in~(c), which belongs to the passage from the
tower to its levels.
\end{outline}

\begin{proof}[Proof of (a), (b), (d), (e), (f) and of the tower in~(c) (not formalized)]
(a) Write $\tilde M=PQ$ with inner dimension $\rho=\irk\tilde M$. The permutation representation is a ring homomorphism, so
$M_X=P_XQ_X$ has rank at most $\card X\rho$ and nullity at least $\card X(N-\rho)$.

(b) Apply~(a) to the bipartite double cover $H\times K_2=\mathcal T/\pi^+$, on which every edge is twisted. With $M=I+B$, its
lift is $\begin{psmallmatrix}I&B\\B&I\end{psmallmatrix}$, and $(x,y)$ lies in its kernel exactly when $x=By$ and $(I+B^2)y=0$;
so its nullity is the nullity of $M^2=I+B^2$, at most~$N$. Hence $2Nr_U(M)\le N$. The bound is also contained in
\emph{Covers and Local Generation}~\cite[Proposition~3.2]{Covers2026}, a lower bound $\lceil N/2\rceil$ for the inner rank
over every field, whose separation hypothesis holds for graph templates and whose proof uses the specialization argument
of~(a).

(c) A double cover of level~$n$ is given by voltages in $\Z/2=\{1,g\}$ on the edges: the subgroup $\Lambda_{n+1}$ of
index~$2$ in~$\Lambda_n$ defines the double cover $\mathcal T/\Lambda_{n+1}\to\mathcal T/\Lambda_n$, and the lift $M_{n+1}$ is
the lift of~$M_n$ along it, with the twisted edges given by the voltages. The step above then gives the recursion at every
level. For symmetric checks, the
adjacency checks $M=I+A$, the recursion is the lift identity of \emph{Covers and Local Generation}~\cite[Lemma~9.1]{Covers2026},
for arbitrary binary lifts of finite connected graphs; the monotonicity of the normalized nullity under refinements with
$2$-group kernels in characteristic~$2$ is its Lemma~4.1, and it is also the index-$2$ case of the argument of Bergeron,
Linnell, Lück and Sauer~\cite[Lemma~4.1 and Theorem~1.6]{BLLS2014} (see also Kionke~\cite[Lemma~2.6]{Kionke2017}). Part~(c)
gives the exact identity, with the twist pairing, for every unit-diagonal, possibly asymmetric, binary check, along arbitrary
towers of index-$2$ subgroups.

(d) If every level were non-bipartite, the sets $\overline{\Lambda_n}\cap(\hat\pi\setminus\hat\pi^+)$ would form a decreasing
sequence of nonempty compact sets, so $C\not\le\hat\pi^+$. Conversely, a bipartite level gives
$C\le\overline{\Lambda_n}\le\hat\pi^+$. Covers of bipartite graphs are bipartite, and every storage code on a bipartite graph has
rate at most~$\frac12$~\cite[Corollary~6.9]{FourthPaper2026}.

(e) A free pro-$2$ group of positive rank maps onto~$\Z_2$ with a free pro-$2$ kernel, so it is free-by-$\Z_2$. By
Jaikin-Zapirain and Souza~\cite[Theorem~8.1]{JaikinSouza2026}, applied to the finitely generated subgroup $\iota(\pi)\cong\pi$,
the normalized dimensions of the kernels converge to the dimension over the universal division ring of~$\F_2[\iota(\pi)]$; in the
case of a finitely generated free pro-$2$ group, which is all that \cref{prop:eqtower} uses, this is
Jaikin-Zapirain's~\cite[Proposition~3.5]{JaikinZapirain2024}. Since $\iota$ is injective, this division ring is Cohn's free field
of $\F_2[\pi]$, and the limit is $r_U(M)$.

(f) The Clebsch graph contains the $4$-cube, whose Gray code is a Hamiltonian cycle, so it has a spanning $2$-factor, and
$I+A$ has unit diagonal and nonzero scalars on all edges. For such checks, \emph{Covers and Local Generation} shows that the
abelianization of~$\tilde M$, over the Laurent polynomial ring in the chord generators, has nonzero
determinant~\cite[Lemma~6.1]{Covers2026}, with a written proof there. A factorization $\tilde M=PQ$ through fewer than $N$
columns would specialize to a factorization of this invertible matrix, so $\irk\tilde M=N=16$, as in the proof
of~\cite[Theorem~6.2]{Covers2026}, and $r_U(M)=0$. As an independent computational check, the script
\leanfile{tools/clebsch_lift.py} constructs an explicit $16$-sheeted lift~$X_0$ of the Clebsch graph, by permutation voltages
on the $25$ edges outside a breadth-first spanning tree, and computes, by exact elimination over~$\F_2$, that the lift of
$I+A$ has nullity~$10$; by~(a), $10\ge16(16-\irk\tilde M)$, so again $\irk\tilde M=16$.
\end{proof}

\begin{proposition}[Towers that never become bipartite]\label{prop:eqtower}
\notformalized
There are towers of normal subgroups of $\pi_1(\mathrm{Clebsch})$ with trivial intersection whose covers are non-bipartite at
every level and have girth tending to infinity, and along which the rate of the lifted Clebsch code tends to~$0$.
\end{proposition}

\begin{informalproof}
The Clebsch graph has $16$ vertices and $40$ edges, so $\pi_1(\mathrm{Clebsch})$ is free of rank~$25$. The kernel
$K=\ker(F_9\to\Z/3)$ of a surjection from the free group of rank~$9$ is free of rank $1+3\cdot8=25$ by Schreier's formula. It is
dense in the pro-$2$ completion $\hat F_9$, since for every normal subgroup $N$ of $2$-power index the index of $KN$ divides both
$3$ and a power of~$2$; and $K$ embeds in $\hat F_9$, since free groups are residually~$2$. A character of~$K$ is continuous for
the induced topology exactly when it is the restriction of a character of~$F_9$, and restriction is injective on characters
because $F_9/K$ has order~$3$; so exactly $2^9$ of the $2^{25}$ characters of~$K$ are continuous. Since $\mathrm{Aut}(F_{25})$
maps onto $\mathrm{GL}_{25}(\F_2)$, it acts transitively on nonzero characters, so there is an isomorphism
$\varphi\colon\pi_1(\mathrm{Clebsch})\to K$ for which $\chi\circ\varphi^{-1}$ is not continuous. For a descending chain of open
normal subgroups $U_j\le\hat F_9$ with trivial intersection, put $\Lambda_j=\varphi^{-1}(K\cap U_j)$. These are normal with
trivial intersection; $\Lambda_j\not\le\pi^+$ for every~$j$, since otherwise $\chi\circ\varphi^{-1}$ would vanish on $K\cap U_j$
and be continuous; and the girth tends to infinity, because only finitely many conjugacy classes have bounded translation
length and each nontrivial element eventually leaves the normal subgroups~$\Lambda_j$. By \cref{thm:eqcrux}(e), with $\iota$ the
composite of $\varphi$ and the inclusion $K\subseteq\hat F_9$, and by \cref{thm:eqcrux}(f), the rate tends to $r_U=0$; the limit
also follows from~\cite[Theorem~6.2]{Covers2026}.
\end{informalproof}

\paragraph{Why it matters.} A construction cannot obtain its coding gain from a fixed small code lifted along towers that
become bipartite, nor, for binary linear codes, along towers pulled back injectively from free pro-$2$ groups: bipartite levels
have rate at most~$\frac12$ for every code (\cref{thm:eqcrux}(d)), lifting a binary check along double covers never raises the
rate (\cref{thm:eqcrux}(c)), and along towers of type~(e) the rate tends to the universal rate, at most~$\frac12$, and to~$0$
for the Clebsch code, whose rate is $\frac58$ on its own graph~\cite[Corollary~6.11]{FourthPaper2026}. The theorem leaves
open: normal
towers with trivial intersection and $C\not\le\hat\pi^+$ that are not of type~(e), such as $2$-adic analytic towers, where a
positive limit for the Clebsch code would contradict the mod-$2$ case of the Lück approximation conjecture in positive
characteristic, which Jaikin-Zapirain and Souza discuss for chains with trivial intersection~\cite[Section~8]{JaikinSouza2026};
non-normal towers; towers with short even cycles at positive density; and codes that are not lifts of a code on the base.
Trivial intersection matters: for the Clebsch adjacency check, the cyclic covers of degree $2^j$ with the voltage on a single
edge form a normal tower of binary lifts whose intersection is not trivial; odd cycles that avoid that edge lift to odd
cycles, so no level is bipartite, and the dimension $N/2+2$ at every level~\cite[Theorem~9.5]{Covers2026} makes the rate tend
to~$\frac12$, not to $r_U=0$.

\paragraph{Relation to earlier work.} \emph{Covers and Local Generation in Linear Storage Codes}~\cite{Covers2026} contains
several parts of the theorem. Its lift identity~\cite[Lemma~9.1]{Covers2026} is the recursion of~(c) for symmetric binary
checks, and its Lemma~4.1 gives the monotonicity in~(c) for towers of normal subgroups. Its Proposition~3.2 contains~(b). Its Lemma~6.1 and Theorem~6.2, full inner rank and rate tending
to~$0$ along every chain from a faithful embedding into a free-by-$\Z_p$ pro-$p$ group, for every check with nonzero edge
entries on a graph with a spanning $2$-factor, give~(f) without computation and the limit in \cref{prop:eqtower}. It also
proves, under pro-$p$ and separation hypotheses, an asymptotic half-rate bound for homogeneous checks on a fixed free-group
voltage template in every characteristic~$p$, uniformly over coefficient fields and values~\cite[Theorem~4.3]{Covers2026},
through Jaikin-Zapirain and Souza's Theorem~8.1, and the exact finite bound $N/2+2$ for the Clebsch adjacency code along
every sequence of connected binary lifts~\cite[Theorem~9.2]{Covers2026}. New here are the twist-pairing recursion of~(c) for
asymmetric binary checks along arbitrary towers of index-$2$ subgroups, part~(d), and the towers of \cref{prop:eqtower}, which
never become bipartite, have girth tending to infinity and are pulled back injectively from~$\hat F_9$.

\subsection{Scalar codes at degree four}\label{sec:crux-deg4}

Let $F$ be a field and $M\in F^{V\times V}$ a matrix with unit diagonal, $N=\card V$. The \emph{recovery set} of~$v$ is
$S_v=\{u\ne v:M_{vu}\ne0\}$, and we say that $u$ \emph{points to}~$v$ if $v\in S_u$. The matrix has \emph{no recovery triangle}
if $q\notin S_p$ whenever $p,q\in S_r$ and $p\ne q$; this holds whenever the off-diagonal support of~$M$ lies in the edges of a
triangle-free graph. Following the fourth paper's Theorem~S~\cite[Theorem~6.13]{FourthPaper2026}, consider at a vertex~$v$ the
conditions
\begin{itemize}[leftmargin=3em]
\item[$(\star\star)$] $\card{S_v}\le2$, or $\card{S_v}=3$ and $r\in S_v$ whenever $v\in S_r$;
\item[$(\star4)$] $\card{S_v}\le3$, or $\card{S_v}=4$ and $r\in S_v$ whenever $v\in S_r$.
\end{itemize}
If the support of~$M$ lies in a graph of maximum degree at most~$4$, every vertex satisfies $(\star4)$, since
$\card{S_v}=4$ forces $S_v$ to be the whole neighbourhood. A vertex violating $(\star\star)$ is \emph{heavy}, and $H$ is the set
of heavy vertices. The \emph{separation} conditions are: (S1) no vertex has two heavy vertices in its recovery set; (S2) no
heavy vertex has a heavy vertex in its recovery set. If the support lies in a graph~$G$ and the heavy vertices are pairwise at
distance at least~$3$ in~$G$, both hold.

\begin{theorem}[Theorem SH: separated heavy vertices]\label{thm:SH}
\formal{TheoremSH.lean}{theoremSH, theoremSH_graph}
Let $M$ have unit diagonal and no recovery triangle, and let every vertex satisfy $(\star4)$. If $\card H\le2$, or if the heavy
vertices satisfy (S1) and~(S2), then $\dim\ker M\le\lfloor N/2\rfloor$, over every field. In graph form: if $G$ is triangle-free
of maximum degree at most~$4$, the off-diagonal support of~$M$ lies in the edges of~$G$, and the vertices violating
$(\star\star)$ number at most two or lie pairwise at distance at least~$3$ in~$G$, then $\rank M\ge\lceil N/2\rceil$.
\end{theorem}

Theorem~S of the fourth paper is the case $H=\emptyset$. Since $\dim\ker M^\T=\dim\ker M$ and the absence of recovery
triangles is invariant under transposition (\lean{finrank_ker_transpose}, \lean{recTriangleFree_transpose}), the same holds
when $M^\T$ satisfies $(\star4)$ and its heavy vertices, defined through the sets $\{u\ne v:M_{uv}\ne0\}$, satisfy the
hypothesis instead (\lean{theoremSH_transpose}).

\begin{corollary}[Barrier and matching forms]\label{cor:SH}
\formal{TheoremSH.lean}{finrank_ker_le_barrier_SH, rho3, finrank_ker_le_matching_rho3, finrank_ker_le_matching_rho3_transpose}
Let $M$ have unit diagonal and no recovery triangle, let $U\subseteq V$, and let $V\setminus U$ be partitioned into blocks with
$M_{uv}=0$ for $u,v$ in different blocks. If on every block~$B$ every vertex satisfies $(\star4)$ on~$B$ and the vertices
violating $(\star\star)$ on~$B$ are at most two or satisfy (S1) and~(S2) on~$B$, then
$\dim\ker M\le\card U+\sum_B\lfloor\card B/2\rfloor$. If the support of~$M$ lies in a triangle-free graph~$G$ of maximum degree at
most~$4$, then
\[
  \dim\ker M\;\le\;\nu(G)+\card H-\max\bigl(\rho_3(H),\min(\card H,2)\bigr),
\]
where $\nu(G)$ is the matching number and $\rho_3(H)$ the largest size of a subset of~$H$ at pairwise distance at least~$3$
in~$G$; the same holds with the heavy set of~$M^\T$ in place of~$H$.
\end{corollary}

\begin{outline}
In Lean the unit diagonal and the absence of recovery triangles form the hypothesis \lean{MotionDepthCodes.ThmSHyp}, and
the other hypotheses of \cref{thm:SH}, relative to a vertex set~$W$, form \lean{SHOn}: $(\star4)$ on~$W$, and at most two heavy
vertices or (S1) and (S2) on~$W$, where the heavy vertices are those that violate $(\star\star)$ on~$W$ (\lean{Star4On},
\lean{heavySet}, \lean{StarOn}, \lean{S1On}, \lean{S2On}). The half bound holds for the code of~$M$ on every such~$W$
(\lean{finrank_subCode_le_half_SH}). The kernel of $M$ meets $\{x_U=0\}$ in a
subspace of codimension at most $\card U$, which lies in the code on $V\setminus U$; this code splits over the blocks
(\lean{MotionDepthCodes.finrank_subCode_biUnion_le}, for the graph of the nonzero entries), and the half bound applies to each
block (\lean{finrank_ker_le_barrier_SH}). For the matching form (\lean{finrank_ker_add_card_le_matching}), let $D$ be a set of
heavy vertices with at most two elements or at pairwise distance at least~$3$, that is, not adjacent and without a common
neighbour (\lean{Far3}), and put $U=H\setminus D$. The Tutte--Berge formula~\cite{Berge1958}, in the barrier form of
lean-machine (\lean{SimpleGraph.exists_card_add_sum_card_div_two_eq_matchingNumber}), gives a set~$U'$ with
$\card{U'}+\sum_B\lfloor\card B/2\rfloor=\nu(G)$ over the components~$B$ of $G-U'$. Apply the barrier form with the set
$U\cup U'$ and the blocks $B\setminus U$. They are separated, and on each of them the heavy vertices lie in~$D$, since a vertex
heavy on a subset is heavy (\lean{heavySet_mono}); so the hypotheses hold on each block, (S1) and (S2) following from the
distances. Hence $\dim\ker M\le\card U+\card{U'}+\sum_B\lfloor\card{B\setminus U}/2\rfloor\le\card H-\card D+\nu(G)$. A largest
set at pairwise distance at least~$3$ (\lean{rho3}, \lean{exists_rho3}) and a set of $\min(\card H,2)$ heavy vertices are such
sets~$D$, which gives the bound as stated (\lean{finrank_ker_le_matching_rho3}). For $M^\T$, whose support lies in~$G$ as
well, apply the bound to~$M^\T$ (\lean{finrank_ker_le_matching_rho3_transpose}).
\end{outline}

\paragraph{Why it matters.} The theorem extends the fourth paper's Theorem~S, the case $H=\emptyset$, to matrices with few
or scattered heavy vertices at degree four. A scalar linear code $\ker M$ beating $\lfloor N/2\rfloor$ by $\eps N$ on a
triangle-free graph of maximum degree four needs $\card H-\rho_3(H)\ge\eps N$, for~$M$ and for~$M^\T$ alike: a linear number
of heavy vertices that cannot be scattered, that is, dense clusters of vertices whose rows recover them from four neighbours
or from three non-reciprocally. For all codes, guessing numbers of triangle-free graphs need not equal the lower bound
given by fractional clique covers, as Cameron, Dang and Riis showed with the Higman--Sims graph~\cite{CameronDangRiis2016}. This is the regime of the fourth paper's tree code, in which every
interior vertex is heavy and no locally confined argument can work~\cite[Proposition~6.28]{FourthPaper2026}. For linear machines,
\cref{thm:Dlin} needs the linear crux, and the theorem says where a linear gain at degree four would have to live.

\paragraph{Proof of \cref{thm:SH}.} Write $T_v=\{v\}\cup S_v$ for the support of the row $z_v$ of~$M$. The proof works with the
codes of the fourth paper's formalization on vertex sets (\lean{MotionDepthCodes.subCode}): the code on~$W$ consists of the
vectors supported on~$W$ that satisfy the rows of~$W$, and it is the kernel of $M[W,W]$. For a vertex set~$W$, ``on~$W$'' means
that $S_v$ is replaced by $S_v\cap W$ (\lean{rs}) and only rows $r\in W$ are considered. The unit diagonal and the absence of recovery
triangles are hypotheses on all of~$M$ (\lean{MotionDepthCodes.ThmSHyp}). The class~$\mathcal H$ of vertex sets~$W$ on which
$(\star4)$ and the condition on the heavy vertices hold (\lean{SHOn}) is \emph{hereditary} (\lean{SHOn.mono}): $(\star4)$
passes to subsets, since a recovery set either keeps its size or loses elements, and the reciprocity condition only loses
in-neighbours; a vertex heavy on $W'\subseteq W$ is heavy on~$W$, for the same reason
(\lean{heavySet_mono}); so $\card H\le2$, (S1) and (S2) on~$W$ imply them on~$W'$. We prove the half bound for every
$W\in\mathcal H$ by strong induction on~$W$ (\lean{finrank_subCode_le_half_SH}). Let $W\in\mathcal H$ violate it while all its
proper subsets satisfy it (\lean{MinCE}); below, $V$ stands for~$W$, $N=\card V$, $\mathcal C$ is its code and $H$ its set of
heavy vertices. The two cases, $\card H\le2$ and the separated case, are treated together.

For $X\subseteq V$ let $c(X)$ be the dimension of the codewords vanishing on~$X$ (\lean{cdim}). For $u\in X$, the
\emph{pending set} of~$u$ at~$X$ is $P_X(u)=S_u\setminus X$ (\lean{pend}), and $u$ is \emph{closed} at~$X$ if
$P_X(u)=\emptyset$. A \emph{move} at~$X$ is a
pair $(r,P_X(r))$ with $r\in X$ and $P_X(r)\ne\emptyset$; it leads to $X\cup P_X(r)$, and $P_X(r)$ is its \emph{block}.

\begin{lemma}[Axioms of a minimal counterexample]\label{lem:SH-axioms}
\formal{SHBase.lean}{MinCE.odd, MinCE.dim, MinCE.ub, MinCE.two_le_rs, MinCE.out_card}
\begin{enumerate}[label=(\alph*)]
\item $N=2m+1$ and $\dim\mathcal C=m+1$.
\item For every nonempty $X\subseteq V$, $c(X)\le\lfloor(N-\card X)/2\rfloor$.
\item $\card{S_v}\ge2$ for every~$v$.
\item For every nonempty $X\subsetneq V$, the set $\mathrm{Out}(X)=\{v\in X:S_v\not\subseteq X\}$ has at least two elements if
$\card X$ is even and at least one if $\card X$ is odd.
\end{enumerate}
\end{lemma}

By~(a), $\rank M=m$. By duality, (b) says that $\dim(\mathcal C^\perp\cap F^X)\le\lceil\card X/2\rceil-1$, where
$\mathcal C^\perp$ is the row space; in particular the rows supported in~$X$ span at most $\lceil\card X/2\rceil-1$
dimensions.

\begin{outline}
(a) Minimality gives $\dim\ker M[Y,Y]\le\lfloor\card Y/2\rfloor$ for $Y\subsetneq V$. The codewords vanishing at a vertex~$v$
lie in the code on $V-v$, so $\dim\mathcal C\le\lfloor(N-1)/2\rfloor+1$, which exceeds $\lfloor N/2\rfloor$ only if $N=2m+1$ and
$\dim\mathcal C=m+1$ (\lean{MotionDepthCodes.start_facts}); then $c(\{v\})=m$ for every~$v$ (\lean{MinCE.single}). (b) The
codewords vanishing on~$X$ lie in the code on the proper subset $V\setminus X$
(\lean{MotionDepthCodes.subCode_inf_zeroOn_le}). (c) If $S_v=\emptyset$, the row of~$v$ forces $x_v=0$, so $c(\{v\})=m+1$. If
$S_v=\{u\}$, the codewords vanishing at~$v$ vanish at~$u$, so $c(\{v,u\})=c(\{v\})=m>\lfloor(N-2)/2\rfloor$
(\lean{cdim_move}, \lean{MinCE.single_odd_false}). Both contradict~(b). (d) Let $I=X\setminus\mathrm{Out}(X)$ and let $Q$ be
the space of the codewords vanishing on $\mathrm{Out}(X)$, of dimension at least $m+1-\card{\mathrm{Out}(X)}$. The rows of the
vertices of~$I$ are supported in~$X$, so restriction to~$I$ maps~$Q$ into the code on~$I$, of dimension at most
$\lfloor\card I/2\rfloor$, and its kernel is the space of the codewords vanishing on~$X$, of dimension at most
$\lfloor(N-\card X)/2\rfloor$ (\lean{exists_restrict}). So
$m+1-\card{\mathrm{Out}(X)}\le\lfloor\card I/2\rfloor+\lfloor(2m+1-\card X)/2\rfloor$, which gives the claim.
\end{outline}

\emph{Good states and runs.} A nonempty set~$X$ is a \emph{good state} if equality holds in~(b) (\lean{Good}). The row of
the mover makes one pending coordinate a combination of the others, so a move lowers~$c$ by at most $\card{P_X(r)}-1$
(\lean{cdim_move}). Together with \cref{lem:SH-axioms}(b) this gives the effect of a move from a good state~$X$: a block of
size~$1$ is a contradiction if $\card X$ is odd and leads to a good state if $\card X$ is even; a block of size~$2$ leads to a
good state; a block of size~$3$ leads to a good state if $\card X$ is odd; larger blocks, and blocks of size~$3$ at even
states, may leave the good states. The moves of the three kinds that lead to good states, a block of size~$2$, of size~$1$ at
an even state, or of size~$3$ at an odd state, are \emph{neutral}; they lower~$c$ by exactly $\card P-1$
(\lean{MinCE.good_move}, \lean{MinCE.single_odd_false}). A \emph{run} from~$v_0$ is a sequence of neutral moves starting at
the good state $\{v_0\}$; its first move is $(v_0,S_{v_0})$. A \emph{state} of a run is a set that it reaches after at least
one move. The states are good, contain $v_0$ and~$S_{v_0}$, and every other vertex~$u$ of a state has a \emph{parent}, a
vertex of the state that points to~$u$ (\lean{St}, \lean{St.first}, \lean{St.move}). A heavy vertex~$g$ is \emph{non-closed} at~$X$ if $g\notin X$ or
$P_X(g)\ne\emptyset$; the set $\mathrm{NC}(X)$ of these vertices only shrinks as $X$ grows (\lean{NC}, \lean{NC_anti}).

\emph{Test vectors.} At a neutral move $(r,P)$ from~$X$, the codewords vanishing on~$X$ restrict onto the whole hyperplane of
the row of~$r$ on~$P$: for distinct $u,u'\in P$, some codeword vanishes on~$X$ and on $P\setminus\{u,u'\}$ but not at~$u$, nor,
by the row of~$r$, at~$u'$ (\lean{test_vec}). Indeed, imposing the zeros on $P\setminus\{u,u'\}$ lowers~$c$ by at most
$\card P-2$, while a codeword that also vanishes at~$u$ vanishes at~$u'$ and on all of $X\cup P$, where $c$ is lower by
$\card P-1$. Test vectors and rows replace the arguments with the span of the rows of a run.

\begin{lemma}[Pending sizes, immediate closure, a Hall count]\label{lem:SH-tools}
\formal{SHBase.lean}{St.pend_start, St.pend_light, St.pend_le_three, imm_closure, MinCE.hall_count}
\begin{enumerate}[label=(\alph*)]
\item At every state~$X$ of a run from~$v_0$, $v_0$ is closed, a light vertex $u\in X$ other than~$v_0$ has
$\card{P_X(u)}\le2$, and a heavy vertex $u\in X$ other than~$v_0$ has $\card{P_X(u)}\le3$.
\item (\emph{Immediate closure}.) After a neutral move with a block~$P$ of size at least~$2$, every member~$w$ of~$P$ has
$P_{X\cup P}(w)\ne\emptyset$. Consequently a run cannot end at~$V$ with a last block of size at least~$2$.
\item If $I\subseteq V$ is nonempty and no vertex of~$I$ points to another, then $\card{I\cup\bigcup_{v\in I}S_v}\ge2\card I+1$.
\end{enumerate}
\end{lemma}

\begin{outline}
(a) If $\card{S_u}=3$ and $u$ is light, $(\star\star)$ puts its parent, which points to~$u$, into~$S_u$; if $\card{S_u}=4$,
$(\star4)$ does the same; the parent lies in~$X$. (b) Let $u'\in P$ with $u'\ne w$. If $u'\in S_w$, then $w,u'\in S_r$ form a
recovery triangle. Otherwise a test vector for $(w,u')$ vanishes on~$X$ and on $P\setminus\{w,u'\}$ but not at~$w$; if
$S_w\subseteq X\cup P$, it vanishes on all of~$S_w$, and the row of~$w$ forces it to vanish at~$w$. At~$V$ every vertex is
closed (\lean{rs_subset}). (c) The rows $z_v$, $v\in I$, are supported in $U=I\cup\bigcup_{v\in I}S_v$, and they vanish on $I\setminus\{v\}$ since no
vertex of~$I$ points to another. So the codewords vanishing on $U\setminus I$ vanish on~$U$, and by \cref{lem:SH-axioms}(a) and~(b)
$m+1\le\card{U\setminus I}+c(U)\le\card U-\card I+\lfloor(2m+1-\card U)/2\rfloor$, which gives the claim.
\end{outline}

\begin{lemma}[The odd-state lemma]\label{lem:SH-OS}
\formal{SHPhase.lean}{OSAt, odd_state}
Assume $\card H\le2$, or (S1) and~(S2). For every $t\ge0$, no state $X\ne V$ of a run, with $\card X$ odd and
$\card{\mathrm{NC}(X)}\le t$, exists.
\end{lemma}

\begin{outline}
In Lean the states are the sets with the properties of \lean{St}, which include the states of runs. By induction on~$t$; the
cases that invoke the hypothesis for $t-1$ do not arise when $t=0$. Extend the run by neutral moves only, so that all states are
good.

\emph{Phase O.} At the current odd state $X'\ne V$: (O1) a vertex of~$X'$ with exactly one pending vertex gives a
contradiction; (O2) otherwise, if a vertex~$u$ of~$X'$ has $\card{P(u)}=2$, make that move, which keeps the state odd. If it
yields~$V$, its last block is a pair, contradicting immediate closure. Otherwise, if $u$ is light, continue Phase~O at the new
state; if $u$ is heavy, it is closed at the new state, which therefore has at most $t-1$ non-closed heavy vertices, and the
induction hypothesis gives a contradiction. (O3) Otherwise every vertex of $\mathrm{Out}(X')$ is heavy with exactly three
pending vertices, by \cref{lem:SH-tools}(a) and because $v_0$ is closed. By
\cref{lem:SH-axioms}(d), $\mathrm{Out}(X')\ne\emptyset$, so $t\ge1$; go to Phase~E with $Z=X'$ and some $h\in\mathrm{Out}(Z)$.
Every move enlarges the state, so Phase~O ends.

\emph{Phase E} (\lean{F3}). Let $P=P_Z(h)=\{t_1,t_2,t_3\}$. The move $(h,P)$ is neutral at the odd good state~$Z$; $X_0=Z\cup P$ is an
even state, $h$ is closed, and $\card{\mathrm{NC}(X_0)}\le t-1$. No $t_i$ points to another, since $M$ has no recovery triangle.
If (S1) and (S2) hold, all $t_i$ are light by~(S2); if $\card H\le2$, then $\card{\mathrm{NC}(X_0)}\le1$. Below, a \emph{state} is a state reached
from~$X_0$ by neutral moves, and $Y=V\setminus X$; $\card Y$ is odd at even states.

\emph{(F1) No early singles} (\lean{F1}). At an even state with $\card Y\ge3$, no vertex of the state has exactly one pending vertex:
that move is neutral and leads to an odd state $\ne V$ with at most $t-1$ non-closed heavy vertices, contradicting the induction
hypothesis.

\emph{(F2) Pair coincidence} (\lean{F2}). At an even state~$X$ with $\card Y\ge3$, any two pending sets of size exactly~$2$ of
vertices of~$X$ are equal. We argue by induction on~$\card Y$. Let $u_1,u_2\in X$ with $q_1=P_X(u_1)\ne q_2=P_X(u_2)$ of
size~$2$. If $\card Y=3$, then
$q_1\cup q_2=Y$ and $q_1\cap q_2=\{p\}$. A test vector for the move $(u_1,q_1)$ vanishes on~$X$ but not at~$p$. The other two
vertices of~$Y$ lie with~$p$ in $S_{u_1}$ or in $S_{u_2}$, so they are not in~$S_p$, and $S_p\subseteq X$; the row of~$p$ forces
the test vector to vanish at~$p$. If $\card Y\ge5$ and $\card{q_1\cap q_2}=1$, the move $(u_1,q_1)$ leaves $u_2$ with a single
pending vertex at an even state with $\card Y-2\ge3$, contradicting~(F1). So $q_1\cap q_2=\emptyset$. Let $X_2=X\cup q_2$, after
the move of~$u_2$, and $X_{12}=X\cup q_1\cup q_2$, after the moves of $u_1$ and then~$u_2$; both have last block~$q_2$. Let
$c\in q_2$. By immediate closure at~$X_2$, $P_{X_2}(c)\ne\emptyset$, and by~(F1) at~$X_2$ it does not have size~$1$. If it has
size~$2$, the inductive hypothesis at~$X_2$, applied to $c$ and~$u_1$, whose pending set there is still~$q_1$, gives
$P_{X_2}(c)=q_1$; then $P_{X_{12}}(c)=\emptyset$, contradicting immediate closure at~$X_{12}$. So $\card{P_{X_2}(c)}\ge3$ and $c$
is heavy, by \cref{lem:SH-tools}(a). Then both members of~$q_2$ are heavy, non-closed and in~$S_{u_2}$, contradicting~(S1),
and, if $\card H\le2$, the bound $\card{\mathrm{NC}(X)}\le1$ (\lean{F2Mode}).

\emph{(F3) The final count.} By immediate closure at~$X_0$, every~$t_i$ has $P_{X_0}(t_i)\ne\emptyset$, so a heavy member
of~$P$ is non-closed at~$X_0$; and since no $t_i$ points to another, $S_{t_i}\cap X_0\subseteq Z$ for every~$i$. If $V\setminus X_0=\{w\}$, every~$t_i$ points
to~$w$. A test vector for the move $(h,P)$ and the pair $(t_2,t_3)$ vanishes on~$Z$ and at~$t_1$ but not at~$t_2$; the row
of~$t_1$, with $S_{t_1}\subseteq Z\cup\{w\}$, forces it to vanish at~$w$, and then the row of~$t_2$ forces it to vanish at~$t_2$,
a contradiction. So $\card{V\setminus X_0}\ge3$.

If all~$t_i$ are light, then by~(F1) and \cref{lem:SH-tools}(a) $\card{P_{X_0}(t_i)}=2$, and by~(F2) all three pending pairs
equal one pair~$p$. A light $t_i$ with $\card{S_{t_i}}=3$ has $h\in S_{t_i}$ by $(\star\star)$, since $h$ points to~$t_i$, so
$S_{t_i}=\{h\}\cup p$; otherwise $S_{t_i}=p$. Then $\card{P\cup\bigcup_iS_{t_i}}\le6<7$, contradicting \cref{lem:SH-tools}(c)
with $I=P$.

Otherwise, by~(S2), $\card H\le2$ and $\card{\mathrm{NC}(X_0)}\le1$; let $t_3=g$ be heavy and $t_1,t_2$ light. As above,
$t_1$ and~$t_2$ have the same pending pair $p_1=\{a,b\}$. If $\card{V\setminus X_0}\ge5$, make the move $(t_1,p_1)$, reaching an
even state~$X_1$ with $\card{V\setminus X_1}\ge3$. The vertices $a,b$ are light, since $g$ is the only non-closed heavy vertex.
By immediate closure, (F1) at~$X_1$ and \cref{lem:SH-tools}(a) each of them has a pending pair, the same pair~$p_2$
by~(F2). Both $t_1$ and~$t_2$ point
to~$a$, so if $\card{S_a}=3$, $(\star\star)$ would give $S_a\supseteq\{t_1,t_2\}\cup p_2$, four elements; hence $S_a=p_2$, and
likewise $S_b=p_2$, and \cref{lem:SH-tools}(c) with $I=\{a,b\}$ is contradicted, as $\card{\{a,b\}\cup p_2}=4<5$. If
$\card{V\setminus X_0}=3$, then $c(X_0)=1$, and the state $X_1=X_0\cup p_1$ after the move of~$t_1$ has $c(X_1)=0$. A codeword
that vanishes on~$X_0$ and at~$a$ vanishes at~$b$, by the row of~$t_1$, so some codeword~$y$ vanishes on~$X_0$ with $y_a\ne0$. The
rows of $t_1$ and~$t_2$ give $M_{t_1a}y_a+M_{t_1b}y_b=0=M_{t_2a}y_a+M_{t_2b}y_b$. A test vector~$x$ for the move $(h,P)$ and the
pair $(t_2,g)$ vanishes on~$Z$ and at~$t_1$ but not at~$t_2$, and the same rows give $M_{t_1a}x_a+M_{t_1b}x_b=0$ and
$x_{t_2}+M_{t_2a}x_a+M_{t_2b}x_b=0$. The vectors $(x_a,x_b)$ and $(y_a,y_b)$ are annihilated by the nonzero row
$(M_{t_1a},M_{t_1b})$, so they are proportional; hence $M_{t_2a}x_a+M_{t_2b}x_b=0$ and $x_{t_2}=0$ (\lean{two_rows_zero}), a
contradiction.
\end{outline}

Since $\mathrm{NC}(X)\subseteq H$, \cref{lem:SH-OS} with $t=\card H$ says that in a minimal counterexample no run with at least
one move reaches an odd state other than~$V$.

\begin{outline}[Outline of the formal proof of \cref{thm:SH}]
\emph{The start.} The three cases below make up \lean{MinCE.false_SH}. By \cref{lem:SH-axioms}(c), $\card{S_v}\ge2$ for
all~$v$; light vertices have $\card{S_v}\le3$ and heavy ones $\card{S_v}\in\{3,4\}$.
\emph{Case 1: some $v_0$ has $\card{S_{v_0}}=2$.} The move $(v_0,S_{v_0})$ is neutral and gives the odd state $T_{v_0}$. If
$T_{v_0}=V$, immediate closure gives a contradiction; otherwise \cref{lem:SH-OS} does.
\emph{Case 2: no $\card{S_v}=2$, and some heavy~$h$ has $\card{S_h}=3$.} The move $(h,S_h)$ is neutral at $\card X=1$ and gives
the even state $X_0=T_h$ of a run from~$h$, with $h$ closed. This is Phase~E with $Z=\{h\}$ and $P=S_h$: the $t_i$ are light
by~(S2), or $\card{\mathrm{NC}(X_0)}\le\card H-1\le1$; (F1) holds by \cref{lem:SH-OS} with $t=\card{\mathrm{NC}(X_0)}$, and
(F2) and~(F3) give a contradiction.
\emph{Case 3: every light~$v$ has $\card{S_v}=3$ and every heavy~$v$ has $\card{S_v}=4$.} By $(\star\star)$ and $(\star4)$ every
in-neighbour of~$v$ lies in~$S_v$, for every~$v$, so the relation is symmetric. Its graph~$\Gamma$ is triangle-free, with light
vertices of degree~$3$ and heavy ones of degree~$4$. The vertices of odd degree are even in number and $N$ is odd, so $\card H$
is odd and $H\ne\emptyset$; if $\card H\le2$, then $\card H=1$, where (S1) and (S2) hold. Pick $h\in H$ and $v_0\in S_h$. Then
$v_0$ is light by~(S2), $S_{v_0}=\{h,s_1,s_2\}$, and $s_1,s_2$ are light by~(S1), since $v_0$ points to~$h$. The move
$(v_0,S_{v_0})$ is neutral at $\card X=1$, and $X_0=\{v_0,h,s_1,s_2\}$ is an even state. (F1) and~(F2) hold at the even states
reached from~$X_0$, with the same proofs, using \cref{lem:SH-OS} with $t=\card{\mathrm{NC}(X_0)}$. The vertex~$h$, its four neighbours
and $s_1,s_2$ are seven distinct vertices, since $\Gamma$ is triangle-free. If $N=7$, each neighbour of~$h$ has its two further
neighbours among $s_1,s_2$, so $s_1$ has the four neighbours of~$h$ as neighbours, although it is light. So $N\ge9$. Since
$\Gamma$ is triangle-free, $S_{s_i}\cap X_0=\{v_0\}$, so $P_{X_0}(s_i)$ has size~$2$, and by~(F2) both pending sets equal one
pair~$p_1$, as $\card{V\setminus X_0}\ge5$. Make the move $(s_1,p_1)$; $X_1=X_0\cup p_1$ is an even state with
$\card{V\setminus X_1}\ge3$. By~(S1) some $a\in p_1$ is light. As $a\in S_{s_1}\cap S_{s_2}$, symmetry gives $S_a=\{s_1,s_2,x\}$,
so $P_{X_1}(a)\subseteq\{x\}$; by immediate closure it is nonempty, and $a$ has a single pending vertex, contradicting~(F1).

All cases are contradictory, which proves the half bound on every set of~$\mathcal H$ (\lean{finrank_subCode_le_half_SH}); the
case $W=V$ is \cref{thm:SH} (\lean{theoremSH}). For the graph form (\lean{theoremSH_graph}), a triangle-free support gives no
recovery triangle (\lean{MotionDepthCodes.thmSHyp_of_cliqueFree}), maximum degree at most~$4$ gives $(\star4)$
(\lean{star4_of_degree_le_four}), heavy vertices at pairwise distance at least~$3$ satisfy (S1) and~(S2)
(\lean{sep_of_far3}), and the rank is $N$ minus the nullity.
\end{outline}

\begin{remark}[Where the argument stops]\label{rem:SH-stop}
Within the odd-state lemma and Case~2 of the start, (S1) is used only at the end of~(F2) and (S2) only to make the members of
a triple block light, and there, without them, the argument fails exactly at two configurations: two vertices whose pending
pairs are disjoint and consist of four heavy vertices, and a triple block of a heavy vertex that contains a heavy vertex. Case~3
uses them also to make $v_0$, $s_1$, $s_2$ and a member of~$p_1$ light. By \cref{thm:SH}, a minimal counterexample to the
general degree-four question has heavy vertices at distance at most two, and the contradiction tree of the argument must break
at a cluster of heavy vertices.
\end{remark}

\begin{remark}[Numerical: small cases]\label{rem:SH-small}
Exact searches over~$\F_2$ for all digraphs satisfying the hypotheses of \cref{thm:SH} with at most nine vertices, with a
positive control that allows recovery triangles, found no counterexample; such small cases cannot test the asymptotic
question.
\end{remark}
